%% sigplan-test — two-column proceedings format (LaTeX twin of sigconf-test.typ).
\documentclass[sigplan]{acmart}

\setcopyright{acmlicensed}
\copyrightyear{2018}
\acmYear{2018}
\acmMonth{6}
\acmDOI{XXXXXXX.XXXXXXX}
\acmISBN{978-1-4503-XXXX-X/2018/06}
\acmConference[Conference'17]{Proceedings of ACM Conference}{June 2018}{Washington, DC, USA}
\acmBooktitle{Proceedings of ACM Conference (Conference'17)}

\begin{document}

\title{A Two-Column Conference Sample}

\author{Ben Trovato}
\email{trovato@corporation.com}
\affiliation{%
  \institution{Institute for Clarity in Documentation}
  \city{Dublin}
  \state{Ohio}
  \country{USA}}

\author{Lars Th{\o}rv{\"a}ld}
\email{larst@affiliation.org}
\affiliation{%
  \institution{The Th{\o}rv{\"a}ld Group}
  \city{Hekla}
  \country{Iceland}}

\author{Valerie B\'eranger}
\email{valerie@inria.fr}
\affiliation{%
  \institution{Inria Paris-Rocquencourt}
  \city{Rocquencourt}
  \country{France}}

\begin{abstract}
  A short proceedings document used to compare the sigplan two-column layout
  between the LaTeX and Typst renderings: the full-width title and author grid,
  the first-column permission block, and the section typography.
\end{abstract}

\begin{CCSXML}
<ccs2012>
 <concept>
  <concept_id>10010147.10010169.10010170</concept_id>
  <concept_desc>Computing methodologies~Massively parallel algorithms</concept_desc>
  <concept_significance>500</concept_significance>
 </concept>
 <concept>
  <concept_id>10010147.10010169.10010171</concept_id>
  <concept_desc>Computing methodologies~Concurrent algorithms</concept_desc>
  <concept_significance>300</concept_significance>
 </concept>
</ccs2012>
\end{CCSXML}

\ccsdesc[500]{Computing methodologies~Massively parallel algorithms}
\ccsdesc[300]{Computing methodologies~Concurrent algorithms}

\keywords{datasets, neural networks, gaze detection, text tagging}

\maketitle

\section{Introduction}
This document exercises the sigconf format's two-column body, the spanning title
block, and the serif-bold Large section headings. Lorem ipsum dolor sit amet,
consectetur adipiscing elit, sed do eiusmod tempor incididunt ut labore et dolore
magnam aliquam quaerat voluptatem. Ut enim aeque doleamus animo, cum corpore
dolemus, fieri tamen permagna accessio potest, si aliquod aeternum et infinitum
impendere malum nobis opinemur. Quod idem licet transferre in voluptatem, ut
postea variari voluptas distinguique possit, augeri amplificarique non possit. At
etiam Athenis, ut e patre audiebam facete et urbane Stoicos irridente, statua est
in quo a nobis philosophia defensa et.

\subsection{Background}
A subsection to check the level-2 heading. Lorem ipsum dolor sit amet,
consectetur adipiscing elit, sed do eiusmod tempor incididunt ut labore et dolore
magnam aliquam quaerat voluptatem. Ut enim aeque doleamus animo, cum corpore
dolemus, fieri tamen permagna accessio potest, si aliquod aeternum et infinitum
impendere malum nobis opinemur. Quod idem licet transferre in voluptatem, ut
postea variari voluptas distinguique possit, augeri amplificarique non possit. At.

\subsubsection{A subsubsection}
A run-in subsubsection heading. Lorem ipsum dolor sit amet, consectetur
adipiscing elit, sed do eiusmod tempor incididunt ut labore et dolore magnam
aliquam quaerat voluptatem. Ut enim aeque doleamus animo, cum corpore dolemus,
fieri tamen permagna accessio potest, si aliquod aeternum et infinitum impendere.

\section{Structures}
The sigplan style overrides several defaults: enumerate labels are
\emph{1.}/\emph{a.}, theorem heads are bold at zero indent with upright notes,
proofs carry an italic unindented label, and captions bold only their label.

\begin{enumerate}
\item First enumerated item, labeled ``1.''\ rather than ``(1)''.
\item Second enumerated item for the label comparison.
\end{enumerate}

\begin{theorem}[Fermat]
A theorem with a note: the head is bold, the note stays upright, and the body is
italic.
\end{theorem}

\begin{definition}
A definition keeps a bold head and an upright body in sigplan.
\end{definition}

\begin{proof}
A proof label is italic and unindented in sigplan, ending with a QED square.
\end{proof}

\begin{figure}[h]
\centering
\rule{4cm}{1cm}
\caption{A sigplan caption whose label is bold while this caption text stays
regular, long enough to wrap onto a second line.}
\end{figure}

\section{Method}
A second top-level section. Lorem ipsum dolor sit amet, consectetur adipiscing
elit, sed do eiusmod tempor incididunt ut labore et dolore magnam aliquam quaerat
voluptatem. Ut enim aeque doleamus animo, cum corpore dolemus, fieri tamen
permagna accessio potest, si aliquod aeternum et infinitum impendere malum nobis
opinemur. Quod idem licet transferre in voluptatem, ut postea variari voluptas
distinguique possit, augeri amplificarique non possit. At etiam Athenis, ut e
patre audiebam facete et urbane Stoicos irridente, statua est in quo a nobis
philosophia defensa et collaudata est, cum id, quod maxime placeat, facere
possimus, omnis voluptas assumenda est, omnis dolor repellendus. Temporibus autem
quibusdam et aut officiis debitis aut rerum necessitatibus saepe eveniet, ut et
voluptates repudiandae sint et molestiae non recusandae. Itaque earum rerum.

\end{document}
